% ---------------------------------------------------------------------------
\section{Properties of the commonality hypergraph}
\label{sec:theory}
% ---------------------------------------------------------------------------

Throughout, $G=(V,E)$ is a finite simple undirected graph, $N[u]=\{u\}\cup N(u)$ is
the closed neighbourhood of $u$, and
\begin{equation}
  J(u,v)=\frac{|N[u]\cap N[v]|}{|N[u]\cup N[v]|}
  \label{eq:jaccard}
\end{equation}
is the commonality of the pair. A set $T=\{i,j,\ell\}$ of three distinct nodes is a
\emph{wedge} if one of its nodes is adjacent to the other two, and it is
\emph{admissible at threshold $t$} if it is a wedge and $J\ge t$ for all three of its
pairs. The admissible triples $\mathcal T_t$ are the hyperedges of a 3-uniform
hypergraph $\mathcal H_t=(V,\mathcal T_t)$. Let $L_t$ be its two-node-overlap graph,
whose vertices are the triples of $\mathcal T_t$, two of them adjacent when they share
exactly two nodes. The communities at threshold $t$ are the node sets of its
components:
\begin{equation}
  \mathcal C_t=\Bigl(\textstyle\bigcup_{T\in K}T\Bigr)_{K \text{ a connected component of } L_t}.
  \label{eq:communities}
\end{equation}
A node belongs to every community whose component contains a triple through it, so
$\mathcal C_t$ is a cover, and nodes in no admissible triple are left unassigned. All
results except Lemma~\ref{lem:ringJ} and Theorem~\ref{thm:window}, which are specific to
Jaccard, hold for any commonality function; proofs are in Appendix~\ref{app:proofs}.

\subsection{Determinism, locality and nesting}

\begin{proposition}[Order independence]
\label{prop:determinism}
$\mathcal C_t$ is a function of $(G,t)$ alone: it does not depend on a processing
order, a random seed, or the order in which triples are discovered.
\end{proposition}
This is the property the original formulation~\cite{deinega2023thesis} lacked: it
deleted each edge once used, so whether a later triple could still be formed depended on
which community had been grown first.

\begin{proposition}[Local query]
\label{prop:local}
Let $T_0\in\mathcal T_t$. Breadth-first search over $L_t$ from $T_0$, in which a
frontier pair $\{x,y\}$ is extended by every $w\in N(x)\cup N(y)$ with
$\{x,y,w\}\in\mathcal T_t$, visits exactly the component of $L_t$ that contains $T_0$,
and evaluates $J$ only on pairs within the closed neighbourhoods of the nodes it
reaches.
\end{proposition}

\begin{theorem}[Nesting]
\label{thm:nesting}
If $t_1\le t_2$, then $L_{t_2}$ is the subgraph of $L_{t_1}$ induced on
$\mathcal T_{t_2}$, and every community of $\mathcal C_{t_2}$ is contained in a
community of $\mathcal C_{t_1}$.
\end{theorem}
Raising $t$ only splits communities or removes nodes; it never merges two of them. The
covers at all thresholds therefore form one hierarchy, which changes only at the
finitely many values $J$ takes. Its two ends carry no information: on a connected graph
the cover at $t=0$ is one community, and at $t=1$ only classes of nodes with identical
closed neighbourhoods remain, so the most accurate threshold is interior
(Proposition~\ref{prop:endpoints} of the Supplemental Material).

\subsection{Exact recovery on overlapping ring lattices}

Consider $r$ rings of sizes $s_i\ge 2k+2$, each node joined to its $k/2$ nearest
neighbours on either side ($k\ge4$ even), and a set of \emph{merges}: a node $v$ of one
ring is absorbed into a node $u$ of another, which inherits $v$'s edges and so belongs
to both ground-truth communities. The merges are \emph{isolated} if each original node
takes part in at most one merge and, on each original ring, any two merge positions are
more than $k$ apart.

\begin{lemma}[Commonality on the lattice]
\label{lem:ringJ}
With isolated merges, let $x,y$ be unmerged nodes at ring distance $d\le k$, $u$ a
merged node at distance $d$ from $x$ along one of its rings, and $z,y'$ neighbours of
$u$ on its two different rings. Then
\begin{equation}
\begin{aligned}
  J(x,y)&=J(d)\equiv\frac{k+1-d}{k+1+d},\\
  J(u,x)&=J_u(d)\equiv\frac{k+1-d}{2k+1+d},\\
  J(z,y')&=\frac{1}{2k+1},
\end{aligned}
  \label{eq:ringJ}
\end{equation}
and every other pair on distinct rings has $J=0$.
\end{lemma}

\begin{theorem}[Exact recovery window]
\label{thm:window}
With at least one merge, all of them isolated, $\mathcal C_t$ equals the ground truth
if and only if
\begin{equation}
\begin{aligned}
  &\frac{1}{2k+1}<t\le t^{+}(k),\\
  &t^{+}(k)\equiv\min\Bigl\{J_u(1),\,\max\bigl(J(3),J_u(2)\bigr)\Bigr\}\\
  &\phantom{t^{+}(k)}=\min\Bigl\{\frac{k}{2k+2},\,\max\Bigl(\frac{k-2}{k+4},\frac{k-1}{2k+3}\Bigr)\Bigr\}.
\end{aligned}
  \label{eq:window}
\end{equation}
Without merges the condition is $0\le t\le J(2)=(k-1)/(k+3)$.
\end{theorem}

Explicitly, $t^{+}(k)=k/(2k+2)$ for $k\ge8$, so the window widens towards $(0,\tfrac12)$
as $k$ grows: $(1/17,4/9]$ for $k=8$ and $(1/33,8/17]$ for $k=16$. Both edges have a
direct reading. The lower edge is the price of overlap: a shared node is the only bridge
between its two rings, and the bridge is cut once $t$ exceeds the commonality $1/(2k+1)$
of the rings' members across it. The upper edge is set by the shared nodes themselves:
their commonality with either ring is roughly halved by their double membership, and
above it they drop out. For one isolated merge between rings of degrees $k_a$ and $k_b$
the same argument gives
\begin{equation}
\begin{aligned}
  &\frac{1}{k_a+k_b+1}<t\le\min\{U(k_a,k_b),U(k_b,k_a)\},\\
  &U(x,y)=\min\Bigl\{\frac{x}{x+y+2},\,\max\Bigl(\frac{x-2}{x+4},\frac{x-1}{x+y+3}\Bigr)\Bigr\},
\end{aligned}
  \label{eq:windowab}
\end{equation}
and a family with several merge types is recovered exactly if and only if $t$ lies in the
intersection of their windows; once $k_a\ge2k_b^2-2$ the windows of dense--sparse and
sparse--sparse merges are disjoint, and no single threshold recovers a family with
both. We verified these windows exactly, in rational arithmetic, for $k=4,\dots,24$ and
eleven degree pairs, together with the smaller-ring and multi-ring cases of
Sec.~\ref{app:proofsmore} of the Supplemental Material.

\begin{proposition}[Shared pairs]
\label{prop:sharedpairs}
Let $A$ and $B$ be ground-truth communities, and let the pair $\{u_1,u_2\}\subseteq
A\cap B$ lie in an admissible triple $T_A\subseteq A$ and in an admissible triple
$T_B\subseteq B$. Then $T_A$ and $T_B$ lie in the same component of $L_t$, whatever the
commonality between $A\setminus B$ and $B\setminus A$.
\end{proposition}
Two communities that share a pair of nearby nodes are thus read as one, much as
$k$-clique percolation reads communities sharing $k-1$ nodes as
one~\cite{palla2005uncovering}. The random merges of our benchmarks are not isolated:
with $o_n=0.05$, between $6\%$ and $61\%$ of the rings of a family are fused with
another just above the lower edge of the window (Table~\ref{tab:doublebonds}).

\begin{corollary}[Pair-disjointness]
\label{cor:pairdisjoint}
At every threshold $t$, all admissible triples that contain a given pair of nodes lie in
one component of $L_t$. In particular, an edge that lies in an admissible triple is
assigned to exactly one community, and the communities of a node are those of the
admitted edges incident to it.
\end{corollary}
No pair of nodes binds two communities, so a node's memberships can be read off its
admitted edges, as in link communities~\cite{ahn2010link,evans2009line}. This shares an
assumption of clique percolation and link clustering that ground truth contradicts in
many social and collaboration networks, where overlaps are more densely connected than
the rest of the communities~\cite{yang2014structure,yang2012community}; in co-purchasing
networks, by contrast, overlaps are small~\cite{yang2014overlapping}.

\subsection{Cost}
\label{sec:cost}

Admissibility is decided for each of the $W=\sum_v\binom{d_v}{2}$ wedges from three pair
commonalities, all counted in one pass over the wedges, so the time and memory are
\begin{equation}
  O(n+m+W)=O\bigl(n[1+\langle d^2\rangle]\bigr)
  \label{eq:cost}
\end{equation}
for $n$ nodes and $m$ edges. For a power-law degree tail of exponent $\gamma$ this is
$O(n)$ when $\gamma>3$ and $O(n+m\log m)$ at $\gamma=3$, the exponent of the
Barab\'asi--Albert model~\cite{albert2002statistical}; for $2<\gamma<3$ it grows as a
power of $n$ set by the largest degree, and how often real tails are that heavy is
debated~\cite{broido2019scale,voitalov2019scale}.

In practice the heavy tail matters less than $W$ suggests. Since
$J(u,v)\le\rho(u,v)\equiv(\min(d_u,d_v)+1)/(\max(d_u,d_v)+1)$, a wedge containing an edge
with $\rho<t_{\min}$ cannot be admitted at any $t\ge t_{\min}$ and can be discarded
before scoring, provided the remaining commonalities are computed from the complete
neighbourhoods (Sec.~\ref{app:proofsmore}). The excess wedges of a heavy tail are exactly
of this kind: they join a hub to neighbours of far smaller degree. On thirty-two SNAP
networks with up to $4\times10^6$ nodes, $W/(m\ln m)$ ranges from $0.15$ on road networks
to $268$ on a web graph with hubs of degree $8.4\times10^4$, but after pruning at
$t_{\min}=0.02$ it is at most $12$ on all of them (Table~\ref{tab:S-wedges}); on YouTube
it falls from $33$ to $3.7$. The cost is therefore that of the $\gamma=3$ case even where
$W$ itself is not. Once the triple scores are stored, every further threshold costs
$O(W)$, so a whole sweep of the threshold costs one scoring, whereas a sweep of Leiden's
resolution or Infomap's Markov time costs a full run per value. Clique percolation is the
costly method: the number of maximal cliques can grow exponentially with
$n$~\cite{moon1965cliques}, and on one of the fourteen smaller networks of
Sec.~\ref{sec:real} it is out of reach.
